\documentclass[11pt,a4paper]{article}
\usepackage[T1]{fontenc}
\usepackage{lmodern,microtype}
\usepackage[a4paper,margin=25mm]{geometry}
\usepackage{amsmath,amssymb,booktabs,array,tabularx}
\usepackage[dvipsnames]{xcolor}
\usepackage[colorlinks=true,linkcolor=MidnightBlue,citecolor=MidnightBlue,urlcolor=MidnightBlue]{hyperref}
\usepackage{fancyhdr}
\pagestyle{fancy}
\fancyhf{}
\fancyhead[L]{\small Rank-five fusion categories}
\fancyhead[R]{\small Review and revision}
\fancyfoot[C]{\thepage}
\setlength{\headheight}{14pt}
\setlength{\emergencystretch}{2em}
\setlength{\parskip}{3pt}
\newcommand{\FPdim}{\operatorname{FPdim}}
\newcommand{\Tr}{\operatorname{Tr}}
\newcommand{\N}{\operatorname{N}}
\newcommand{\issue}[1]{\subsection{#1}}
\hypersetup{pdftitle={Rank-five classification: mathematical review and revision report},pdfauthor={}}
\title{Rank-five classification\\[4pt]\large Mathematical review and revision report}
\author{}
\date{24 November 2026}
\begin{document}
\maketitle

\section{Outcome and scope}

The attached manuscript has been revised throughout, including its proof
organization, statements, references, computational explanation and source
package. The main conclusions are retained:
\[
\begin{gathered}
\text{every complex rank-five fusion category has a pseudo-unitary Galois
conjugate}\\
\text{and admits a spherical structure},\\
\mu\le4,\qquad \FPdim(\mathcal C)\le48+24\sqrt3,\\
\text{exactly seventeen based rings, each with a unitary realization}.
\end{gathered}
\]
The central simplification is to prove Galois pseudo-unitarity directly in
Theorem~A. This was already a conclusion of the original paper; the revision
obtains it from the descent argument itself. The separate pseudo-unitarity
proof and the subsequent signed branches of the final verifier become
unnecessary.

The review found genuine statement and verifier defects, as well as several
places where a necessary hypothesis or implication was left implicit.
These have been repaired. No unrepaired counterexample to the main
conclusions, or fatal gap in their revised dependency chain, was found.
This assessment combines mathematical reading, source checks and exact
computational replay. It is not a formal proof-assistant verification.

The review covered the complete included manuscript, the structural
programs, the bounded ring generator, the certificate logic and the
realization data. The difficult descent, invertible-object and perfect-category
arguments received separate checks. The discussion below distinguishes
corrections to the original work from qualifications needed by newly
consolidated statements.

\section{The proof after simplification}

\issue{First obtain a positive spherical conjugate}

Let \(R=K_0(\mathcal C)\), and let \(T\) be the signed quotient of its
sphericalization. Write \(D=\FPdim(\mathcal C)\) and let \(E\) be the
intrinsic global dimension. The ordinary and signed dimension characters
have codegrees \(D\) and \(E\). Positivity gives
\[
 \sigma(E)\le D
\]
at every relevant embedding. Equality forces the positive lift dimensions
to equal the Frobenius--Perron dimensions. It then forces equality of the
signed and ordinary structure constants in that choice of lifts.
The conjugate category is pseudo-unitary, and its spherical structure
transports back to the original category.

The revised Lemma~3.1 states the repeated central identity once. For the
constituent \(B_\sigma\) indexed by a conjugate of the signed dimension
character,
\[
 \chi(F(B_\sigma))=\frac{f_\chi}{\sigma(f_\chi)},\qquad
 \dim F(B_\sigma)=\frac{E}{\sigma(E)},\qquad
 \FPdim F(B_\sigma)=\frac{D}{\sigma(D)}.
\]
It follows that
\[
 \frac ED\le \sigma\!\left(\frac ED\right).
\]
When the complementary signed codegrees are rational, the corresponding
class is an integral unit. These facts are the common input to the
degree-by-degree arithmetic.

The stronger conclusion is checked at every exit. Nontrivial grading gives
weak integrality and hence pseudo-unitarity. The quadratic and cubic
rational-complement argument gives equality at a suitable embedding.
The exceptional quadratic calculation excludes the alternative.
The quartic stabilizer argument actually gives equality of all positive
dimensions, and the remaining quartic sieves exclude unequal ratios.
The quintic calculation makes the two reciprocal dimensions Galois
conjugate. Thus no separate proof of Galois pseudo-unitarity is needed.

\issue{Then bound the fusion rules}

After Theorem~A, the entire classification part works in a positive
spherical Galois conjugate. This convention is stated once. It preserves
the based ring and all Frobenius--Perron dimensions.

\begin{center}
\begin{tabularx}{\textwidth}{@{}p{.24\textwidth}X@{}}
\toprule
Family & Role in the revised proof\\
\midrule
Weakly integral & The published integral classification gives the integral
case; the dimensional grading handles the nonintegral weakly integral case.\\
Four invertibles & The near-group multiplicity theorems give the five
possibilities directly.\\
Three invertibles, or a fixed involution & The structural arguments force
integrality.\\
An exchanged pair & Norm estimates and one shared Morita construction bound
the remaining quadratic family.\\
Perfect categories & A small tensor generator or the degree of
\(\mathbb Q(D)\) controls the possibilities. The quadratic and quartic
branches are excluded; cubic and quintic branches enter the final bound.\\
\bottomrule
\end{tabularx}
\end{center}

The cubic case admits a particularly short ending. Theorem~15.12 already
proves \(35<D<36\). In this branch all four nonunit dimensions exceed two.
For every nonunit \(X_i\),
\[
 d_i^2=D-1-\sum_{\substack{j\ne i\\j\ne0}}d_j^2<D-13<23.
\]
Hence \(d_i<5\), and the general inequality
\(N_{ij}^{k}\le\min(d_i,d_j,d_k)\) gives the integer bound \(\mu\le4\).
The final census then identifies the cubic ring. The separate reconstruction
of its full character table is unnecessary for this dependency chain.

\issue{Finally classify inside the proved bounds}

The generator constructs all 269 commutative rank-five based rings with
\(\mu\le4\), covering all three duality types. The exact dimension test
retains 114. The stronger Theorem~A allows every exclusion to use necessary
conditions for a pseudo-unitary realization. There are 97 exclusions and
seventeen explicit basis identifications with models.

This last step does not establish unitarity merely by matching two arrays.
The tensor comparison establishes a based-ring isomorphism; the cited
categorical constructions establish existence and unitarity. The revised
paper makes that distinction explicit. It also spells out the Izumi
involution producing the dimension endpoint, rather than leaving its
realization to an opaque model name.

\section{Corrections and qualifications}

\issue{The 4,5,5 exclusion is categorical}

The original final statement of the quadratic \(4,5,5\) argument excluded
a perfect self-dual \emph{ring}. That statement was too strong. The
preceding discussion itself exhibited a based ring at \(J=206\), of
multiplicity \(234\). The contradiction uses central induction,
Frobenius--Schur indicators and categorical Cauchy, so it excludes a
\emph{fusion categorification}, not the abstract based ring.

Proposition~14.7 now states the categorical hypothesis. The analogous
mixed-profile and quartic exclusions have been corrected in the same
way, and the hypothesis that all nonunit dimensions exceed two is
explicit where it is used. The \(J=206\) illustration was removed from
the proof narrative because it is not needed to establish the classification.
Its existence is the reason this correction is substantive rather than
terminological.

\issue{The rank-three theorem requires a visible self-duality reduction}

In the proper-subcategory argument, the invocation of Ostrik's quadratic
rank-three theorem needed its self-duality hypothesis stated.
Proposition~2.2 now first disposes of the non-self-dual rank-three case,
which is pointed by Ostrik's Proposition~3.1. The remaining case is
self-dual, so Theorem~5.1 applies with the correct hypotheses
\cite{Ostrik3}. The proof now cites the original rank-two and rank-three
sources directly.

\issue{The modular normalization must be precise}

The original rational-row argument used the phrase \emph{anomaly one}.
The new common support lemma specifies the needed normalization:
\[
 p^+=\sum_Y\theta_Y\dim(Y)^2=\sqrt{\dim(\mathcal D)}.
\]
Equivalently, the multiplicative central charge is one, and the positively
normalized \(S\) and ordinary twist matrix satisfy
\[
 STS=T^{-1}ST^{-1}.
\]
The weaker convention \(p^+/p^-=1\) does not determine this sign.
For example, the three-fermion pointed modular category has
\(p^+=p^-=-2\), so its multiplicative central charge is \(-1\).
This example illustrates the normalization issue; it is not a
counterexample to any rank-five classification conclusion.

All applications in the manuscript are Drinfeld centers, where the
stronger normalization holds. Lemma~13.4 also explicitly assumes real
rows and an integral rational block idempotent; the support conclusion
requires injectivity on twist-one columns. The observation \(J=J^*\)
for a real row identifies the needed \(S\) and \(S^{-1}\) columns.
These qualifications make the consolidated statement match its proof.

\issue{The inert-prime argument now precedes its uses}

The quadratic profile reduction formerly referred forward to an induction
congruence appearing in the final \(\mathbb Q(\sqrt5)\) discussion, which
in turn used the profile reduction. The congruence was independently
valid, but the organization obscured that independence.

It is now Lemma~14.1. Its hypotheses are a quadratic dimension field in
which three is inert and rational complementary codegrees \(3,t,g\)
with \(3\nmid g\). The local partner argument isolates the two rational
branches. An integral induction functional gives
\[
 a+b\equiv0\pmod3,\qquad 2+a+b+c\equiv0\pmod3,\qquad c\equiv0\pmod3,
\]
a contradiction. The later reductions cite this lemma directly.

\issue{Arithmetic filters and the cyclic cutoff}

In the cyclic quartic parametrization, the two displayed parity conditions
are necessary conditions for integrality; they are not a complete
integrality criterion in every case, notably at \(z=2\). The manuscript
now says so. The enumeration was already overinclusive, which is safe
for an exclusion: the clarification does not remove any candidate.

The cyclic quartic checker contained an initial numerical cap \(251\).
The derived bounds already put every relevant trace below \(95\), so
the cap was harmless, but it concealed the reason the search was finite.
It has been replaced by the proved bound
\[
 U=1+\frac{s}{\min(b_0,b_1)},
\]
with an assertion that both denominator bounds are positive.
The checker now asserts the full output counts \(32,10,0\).

\issue{A real defect in the induction certificate verifier}

The original induction verifier compared a computed obstruction with the
stored obstruction. It could accept a branch when both values were null.
Equality of two null values is not an exclusion certificate.

The revised verifier requires a non-null computed obstruction before
comparing records. A fabricated null-obstruction certificate for the
realized \(\mathrm{Vec}(C_5)\) tensor is rejected, while all ten genuine
induction exclusion certificates still pass. Every distributed exclusion
already had an actual obstruction, so this repair does not change the
classification. Unknown certificate kinds are now rejected explicitly.

\issue{Reference and exposition corrections}

The commutativity citation has been updated from ABDP Proposition~8.2 to
Proposition~9.2 in the current version. The bibliography now records its
journal publication \cite{ABDP}. The near-group replacement cites
Meng--Yu precisely, and the endpoint construction cites Izumi's
Lemma~13.1 and Example~13.10 \cite{MengYu,Izumi}.

Several qualifications were also checked while writing the new
certificate explanation. The residual first moment \(-a_i\) is asserted
only for nonunit simples. The Ramanujan interval for a relative average
of \(\nu_j\) is asserted only when
\(K\cap\mathbb Q(\zeta_j)=\mathbb Q\). These are properties already
respected by the code. They are now visible in the prose.

\section{Substantial simplifications}

\begin{center}
\begin{tabularx}{\textwidth}{@{}p{.27\textwidth}X@{}}
\toprule
Change & Mathematical benefit\\
\midrule
One descent theorem & Galois pseudo-unitarity follows directly; the separate
proof and its unique checker are removed.\\
One central-column lemma & Dimension ratios, the minimal embedding and the
integral-unit conclusion are established once and reused.\\
One support lemma & A prime-local block argument handles rational rows and
the paired real rows used in several later exclusions.\\
One Morita construction & Lemma~11.5 supplies the eight-object dual category
and induction pairings for both the codegree-eleven and codegree-five arguments.\\
Existing integral theorem & Proposition~9.2 reads off \(D\le60,\ \mu\le3\)
from the nine known integral rings instead of reproducing their enumeration.\\
Existing near-group theorem & Corollary~10.1 uses the \(C_4\) and \(C_2^2\)
multiplicity classifications instead of the long \(C_4\) calculation.\\
Cubic dimension estimate & The inequality \(d_i^2<23\) supplies
\(\mu\le4\); the final census identifies the ring.\\
One certificate mode & All exclusions use the positive spherical conjugate.
The signed-table and alternate spherical branches are removed.\\
\bottomrule
\end{tabularx}
\end{center}

The integral simplification introduces an explicit dependence on the
published ABDP classification. The \(C_4\) simplification uses the
Meng--Yu preprint of November 2026, whose theorem applies to arbitrary
complex fusion categories, not merely to unitary realizations.
These dependencies are stated in the revised introduction. The original
claim that the bounds avoid the integral rank-five classification has
therefore been removed.

The common support lemma is a consolidation of the repeated indicator
arguments, with a complete proof and precise hypotheses. Its
Ng--Schauenburg inputs are identified at the level of the rotation,
grouping, cyclotomicity and modular indicator formulas \cite{NS}.
The long quadratic discussion is now a sequence of stated reductions,
with named intermediate conclusions and short proof-stage headings.

The presentation now follows the actual dependency order: descent;
weak integrality and invertibles; shared character tools; perfect
categories by quadratic, cubic and quartic field degree; sharp bounds;
finite classification. The quintic perfect case follows immediately
from the earlier quintic descent analysis. The source has been flattened
into one self-contained \texttt{rank5.tex}, including its bibliography;
obsolete fragments and alternative verifier paths are omitted.

\section{Verification and reproducibility}

\issue{The finite census is unchanged}

\begin{center}
\begin{tabular}{@{}rrrrr@{}}
\toprule
\(\mu\) & Generated & Within dimension bound & Excluded & Realized\\
\midrule
1 & 16 & 16 & 6 & 10\\
2 & 37 & 30 & 27 & 3\\
3 & 82 & 46 & 44 & 2\\
4 & 134 & 22 & 20 & 2\\
\midrule
Total & 269 & 114 & 97 & 17\\
\bottomrule
\end{tabular}
\end{center}

\begin{samepage}
The positive-only certificate classification is:
\begin{center}
\begin{tabular}{@{}lr@{}}
\toprule
Exact exclusion criterion & Count\\
\midrule
Formal-codegree \(d\)-number condition & 76\\
Subcategory dimension divisibility & 2\\
Frobenius--Perron norm obstruction & 3\\
Frobenius--Perron dimension-field obstruction & 2\\
Unit-induction Gram and second-moment conditions & 10\\
Finite higher-indicator trace constraints & 4\\
\midrule
Total excluded & 97\\
\bottomrule
\end{tabular}
\end{center}
\end{samepage}

These are exclusive certificate categories in the distributed data, not
claims that each rejected ring fails only one necessary condition.
The old two-stage grouping of the last nineteen exclusions has been
replaced by their direct positive-spherical reasons.

\issue{What was checked}

The original complete verification succeeded. The revised package runs
eighteen structural programs and then regenerates the bounded census,
checks every retained witness and compares every realized tensor under
all permutations of the four nonunit basis elements.

The four dedicated descent calculations reproduce:
\begin{center}
\begin{tabularx}{\textwidth}{@{}p{.18\textwidth}X@{}}
\toprule
Calculation & Output\\
\midrule
Quintic & 34 ratios, 15,396 trace vectors, 39 signed passes and four joint
passes; each joint pair is Galois conjugate.\\
Quadratic & 331 initial candidates, 50 elementary profiles, 13 Fourier
profiles, 22 dimension tuples and one final tuple; the written unit
argument excludes that tuple.\\
Biquadratic & 2,097 fields, 52 after the field bounds, 30 ratios and no
compatible integral-trace/Fourier data.\\
Cyclic quartic & 105 fields, 221 polynomial profiles, 32 ratio profiles,
ten norm candidates and no compatible Fourier data.\\
\bottomrule
\end{tabularx}
\end{center}

Additional independent checks compared all nine nonempty low-column
enumerations with direct Python enumeration and principal-minor tests,
and all nine residual-column enumerations with the reference Gram
backend. The finite twist-order argument was checked to cover arbitrary
orders, including the joint zero-average case beyond the finite divisor
box and the coherent signs in quadratic relative traces.

The stored \(\operatorname{Ad}(\mathrm{SU}(3)_4)\) tensor was independently
reconstructed from the fifteen alcove weights and fundamental Pieri
multiplication. A compact description of that construction is included
in the model registry. The Izumi dimension-endpoint tensor was checked
against the actual order-two tensor symmetry: inversion on \(C_3\) and
a reflection on the multiplicity space. Its two eigenspaces have
dimension three, giving
\[
 X^2=Y^2=1+Z+3X+3Y,\qquad XY=g+Z+3X+3Y,
\]
and \(D=48+24\sqrt3\).

The final manuscript is compiled with stable theorem references and
citations. Its PDF has been rendered for visual inspection. No
unresolved references, duplicate labels or overflowing text boxes remain.

\issue{How to reproduce the result}

The complete source archive contains the manuscript and review PDFs,
their LaTeX sources, the structural programs, the C++ generators, the
114 witness records, the model tensors and the verification transcript.
With Python 3.12, SymPy, NumPy and a C++17 compiler with OpenMP, run
\begin{verbatim}
python3 verify.py
\end{verbatim}
from the extracted archive directory. Assertions must remain enabled.
The two stages can also be run separately with
\texttt{--bounds-only} and \texttt{--classification-only}.
The checks require no network access or floating-point feasibility solver.
The README gives the three-pass \texttt{pdflatex} build instructions.

\section{Interpretation of the review}

The reductions do not assume the final seventeen-ring list in order to
derive the sharp bounds. They do use the explicitly cited integral and
near-group classifications. The remaining categorical and arithmetic
lemmas are proved before the finite census is used.

Successful programs validate the stated finite enumerations and exact
certificates. Their mathematical force still depends on the categorical
lemmas and the proofs that the search ranges are exhaustive. Those
arguments were reviewed as part of this revision; no claim of a
machine-checked proof of all of fusion-category theory is intended.

The principal remaining burden for a journal referee is concentrated
in the signed central-induction/Fourier argument, the prime-local
indicator limit, the odd-multiplicity Morita construction and the
quartic moment bounds. The revision gives these arguments a single
location, explicit hypotheses and a clear role in the proof. The main
classification and both sharp endpoints are retained.

\begin{thebibliography}{9}
\raggedright
\bibitem{ABDP}
M.~A.~Alekseyev, W.~Bruns, J.~Dong and S.~Palcoux,
\emph{Classifying integral Grothendieck rings up to rank 5 and beyond},
Lett. Math. Phys. \textbf{116} (2026), article~141.
\href{https://doi.org/10.1007/s11005-026-02164-7}{doi:10.1007/s11005-026-02164-7};
\href{https://arxiv.org/abs/2507.07023v2}{arXiv:2507.07023v2}.

\bibitem{Ostrik3}
V.~Ostrik, \emph{Pivotal fusion categories of rank 3},
Mosc. Math. J. \textbf{15} (2015), 373--396.
\href{https://arxiv.org/abs/1309.4822}{arXiv:1309.4822}.

\bibitem{MengYu}
Y.~Meng and Z.~Yu, \emph{The multiplicity bound for near-group fusion
categories of type \(\mathbb Z/4\mathbb Z+k\)},
\href{https://arxiv.org/abs/2611.12191}{arXiv:2611.12191} (2026),
Theorems~1.1 and~4.3.

\bibitem{Izumi}
M.~Izumi, \emph{A Cuntz algebra approach to the classification of
near-group categories}, Proc. Centre Math. Appl. Austral. Nat. Univ.
\textbf{46} (2017), 222--343.
\href{https://arxiv.org/abs/1512.04288v2}{arXiv:1512.04288v2},
Lemma~13.1 and Example~13.10.

\bibitem{NS}
S.-H.~Ng and P.~Schauenburg, \emph{Congruence subgroups and generalized
Frobenius--Schur indicators}, Comm. Math. Phys. \textbf{300} (2010),
1--46. \href{https://arxiv.org/abs/0806.2493}{arXiv:0806.2493},
Lemmas~2.5 and~2.7, Propositions~2.8, 5.7 and~6.1.
\end{thebibliography}
\end{document}
